% Three-tier research draft. Evaluation separates measurements, traces and models.
% Compile this directory independently with pdfLaTeX and BibTeX.
\documentclass[sigplan,10pt,anonymous,nonacm]{acmart}
\setkeys{acmart.cls}{balance=false}
\usepackage{subcaption}
\usepackage{tabularx}
\usepackage{placeins}
\setcopyright{none}
\settopmatter{printacmref=false,printccs=false,printfolios=true}
\acmConference[EuroSys '27]{European Conference on Computer Systems}
 {April 19--23, 2027}{Rabat, Morocco}
\acmYear{2027}
\acmDOI{}
\acmISBN{}
\title{AegisKV: Workflow-Aware KV Checkpointing\texorpdfstring{\\}{ }Across HBM, DRAM, and Shared Storage}
\setlength{\emergencystretch}{1.5em}
\renewcommand{\bibfont}{\normalfont\normalsize}
\newcounter{designalgorithm}
\captionsetup[subfigure]{font=normalsize,labelfont=normalfont,textfont=normalfont}
\AtBeginDocument{\typeout{AEGISKV-GEOMETRY: width=\the\textwidth;
 height=\the\textheight; columnsep=\the\columnsep; baseline=\the\baselineskip}}
\begin{document}
\fancypagestyle{aegispreview}{\fancyhf{}\renewcommand{\headrulewidth}{0pt}
 \renewcommand{\footrulewidth}{0pt}\fancyfoot[C]{\normalsize\thepage}}
\pagestyle{aegispreview}
\begin{abstract}
Multi-agent applications accumulate reusable context while executing
interdependent LLM requests. A worker failure can lose this state and delay
several branches at once, yet protecting every context in fast memory is
expensive. We propose AegisKV, a workflow-aware manager for KV checkpoints
across Peer HBM, shared DRAM and shared storage. Exclusive managed objects
across the three tiers cover a committed checkpoint. AegisKV combines
normal reuse, guarded workflow slack and causal sparse demand to set
protection priorities. Fixed layer/block quotas accelerate selected
requests on a compatible takeover DP; lower-tier renewable credit and
bounded lazy updates manage capacity pressure.
Versioned transitions preserve recovery validity during migration and
explicit fallback. A per-DP Adapter changes total Peer capacity while the
Master updates request membership and hot-block identities without
changing the calibrated quotas or native attention semantics. We connect
physical reconstruction microbenchmarks and held-out trace diagnostics
to reproducible analytical controls and a matched experimental protocol.
In 36 declared Peer-disabled scenarios, the placement model improves on
LRU in 17 and ties in 19. Matched analytical schedules illustrate Peer's
incremental overlap benefit; neither model measures a three-tier system. Historical outage timing drives a separate recovery-load model in which
a lower mean coexists with worse tail latency. Historical analytical
controls expose placement--overlap interactions and a rental regression;
the protocol tests how contention, prediction error, and commit lag affect
the current controller. End-to-end fault injection and normal
service measurements remain necessary to establish the proposed system's
performance.
\end{abstract}

\maketitle
\thispagestyle{aegispreview}
\raggedbottom

\section{Introduction}
\label{sec:introduction}

A multi-agent application completes when its required results are ready:
a tested patch, a report, or an answer assembled from parallel searches.
An LLM request's value therefore depends on its workflow position. A
critical-branch delay postpones completion; another branch may absorb it
before a join. Application-aware serving exploits this structure during
normal execution~\cite{lin2024parrot,pan2025kvflow,ni2026topas}. Failures
raise a further question: which histories should remain quickly
recoverable when one worker interrupts several requests together?

Long contexts make that choice expensive. Agents accumulate reasoning,
code, and tool results; losing their GPU-resident KV can require substantial
reconstruction. Checkpointing avoids that work~\cite{strati2024dejavu,jayakody2026ghostserve},
and load-aware recovery distributes interrupted requests~\cite{cao2026lumen}.
Hierarchical caches provide DRAM and storage capacity~\cite{qin2025mooncake,xie2026strata},
but keeping every history in the faster tier is costly.

Recency and recovery latency alone do not determine placement. A long,
recent history may belong to a branch that tolerates storage access,
while a smaller history lies on the critical path. High slack does not
justify deleting reusable history: subsequent normal calls may repeatedly
reconstruct it. Common-worker failures also couple requests through the
DAG and shared recovery resources. Accelerating only one of two delayed
branches may leave their join unchanged. Finally, demotion is asynchronous:
freeing DRAM before its storage replacement commits invalidates protection.

AegisKV manages checkpoints across Peer HBM, shared DRAM and shared storage.
Its Tracker estimates remaining workflow time; the Master combines
normal reuse, guarded slack and causal block demand. The three managed
tiers hold exclusive objects whose union covers a committed checkpoint.
Calibrated layer, block and checkpoint quotas are fixed per compatible
model/layout. Greedy Peer admission protects zero-slack requests, while
lower-tier renewable credit uses bounded lazy refresh. The Router binds
a compatible takeover DP; the Adapter changes total Peer capacity, and
the Connector executes transfers and reports their outcome. Versioned
manifests, reservations and reader leases preserve recovery validity
during migration, pressure relief and explicit fallback.

Our proposed contributions are workflow- and access-aware protection
priorities; a bounded three-tier controller with safe ownership
transitions; and a matched evaluation protocol grounded in physical
path measurements, held-out trace diagnostics and analytical controls.
The available evidence exposes gains and ties in a historical allocation
model, modest average prediction gains, and queue-tail and rental
regressions. It constrains the design without validating the current
controller. The central empirical hypothesis is that workflow-aware
protection improves finite-resource allocation after accounting for
healthy-service and migration costs.

\section{Background} \label{sec:background} 

This section describes how DAG-based agent workflows are mapped to an LLM serving system and how the system recovers an agent's KV state after a serving-worker failure.

\subsection{DAG-Based Agent Workflow Serving} \label{sec:background:serving} 

A multi-agent application coordinates multiple role-specialized agents to complete a user task. An agent may generate intermediate results, invoke external tools, delegate subtasks, and consume outputs produced by other agents. These dependencies form a directed acyclic graph (DAG), in which each node represents an agent invocation and each edge represents a data or control dependency. 

A \textbf{workflow} runtime maintains the DAG and releases an invocation after its dependencies have been satisfied. Forks allow multiple invocations to execute in parallel, while joins delay downstream invocations until all required predecessors have completed~\cite{lin2024parrot}. 
%
An \textbf{agent} can be invoked multiple times during one workflow. On each {turn}, its input may include an agent-specific instruction, previous model outputs, tool results, and outputs received from other agents. These inputs form an evolving agent history retained by the workflow runtime. 

Figure~\ref{fig:serving_model} shows how agent invocations are mapped to an LLM serving system. Ready invocations are submitted as inference requests and scheduled across model workers. A \textbf{worker} contains one model replica, or a distributed replica spanning a tensor-parallel~\cite{shoeybi2019megatron} group, and may concurrently serve agents from multiple workflows. 
%
Processing a request consists of prefill and decoding. During prefill, the model processes the input tokens and generates a key-value (KV) cache for each transformer layer. During decoding, newly generated tokens append entries to this cache. When a later invocation reuses a prefix of the same agent history, the corresponding KV entries allow the worker to process only the newly appended input rather than prefill the complete history again~\cite{gao2024cachedattention,zheng2024sglang}. Consequently, an active agent accumulates a KV history across its iterative turns. 

The KV cache is distributed according to the model-parallel~\cite{kwon2023vllm} configuration. Under tensor parallelism, each rank stores its local KV shards together with the metadata required by the attention implementation. Reusing a KV cache therefore requires compatibility in the model version, parallel configuration, numerical format, and KV layout. 


\begin{figure}[t]
  \centering
  \includegraphics[width=\columnwidth]{figures/serving_model.pdf}
  \caption{DAG-based agent workflow serving. Numbered agent invocations
    are requests; the runtime retains accepted histories and submits
    ready requests to a router for dispatch across workers. Each worker is one TP=2 replica.
    The illustrative MHA layer has eight Q/KV heads: ranks retain the
    same token positions for different heads. B1--B2 shows prefix reuse;
    hatched suffixes are new tokens. Local attention uses matching Q/KV
    heads, with output-projection contributions summed across ranks.
    Dashed arrows return accepted outputs. Bar lengths and placement
    are schematic; reuse requires matching, available KV.}
  \label{fig:serving_model}
  \Description{Two workflow DAGs sit above accepted token histories and
    tool results. Workflow 1 forks from A1 into B1 followed by B2, and
    into C1, then joins at A2. Workflow 2 forks from D1 into E1 and F1, then joins
    at D2. The central Router sends ready B2 and E1 requests to worker 0,
    and C1 and F1 to worker 1 along separate symmetric paths.
    Each worker contains two ranks. Rank 0 owns Q/KV heads 0 through 3
    and rank 1 owns heads 4 through 7 in an illustrative eight-head MHA
    layer. Each rank shows separate K and V bars for both assigned agents.
    For a given agent the two ranks cover the same token positions with
    different heads. Agent B's bars share a retained prefix and a hatched
    newly appended suffix. Dashed paths return outputs from both workers
    to the runtime. Internal computation is omitted from the drawing.}
\end{figure}


\subsection{Worker Failure and KV Recovery} \label{sec:background:recovery} 


We consider fail-stop failures that make a model worker, consisting of one model replica or a tensor-parallel group, unavailable. Such failures may originate from an accelerator, serving process, host, or communication fault~\cite{cao2026lumen,jayakody2026ghostserve}. The protected failure domain is deployment-specific and may correspond to a process, a distributed model replica, or an entire host. A recovery copy provides protection only if it remains reachable after the corresponding domain fails. For example, host-local DRAM or SSD may survive a process failure but does not protect against a host failure. 

A worker may hold KV states from multiple agents and workflows. Consequently, one worker failure can simultaneously interrupt several branches of the same workflow, as well as agents from different workflows. We treat all invocations whose KV states become unavailable under the same worker failure as \textit{one failure event}. 

We assume that the workflow DAG, accepted tokens, and completed tool results are retained outside the failed worker. Affected invocations can therefore be reissued to a compatible surviving replica with the same model version, parallel configuration, numerical format, and KV layout. We focus on model-worker failures because they remove large, worker-local KV state while leaving the workflow recoverable from its durable inputs. Failures of the workflow runtime or durable token store are orthogonal and can be handled using conventional replication mechanisms. Unlike these control and durable states, KV caches grow with agent histories, making full replication costly and creating the resource-allocation problem studied in this paper. We assume that at least one advertised recovery domain remains available. 

Without a surviving KV copy, the replacement replica must prefill each affected agent's retained history before generation can resume~\cite{strati2024dejavu}. Alternatively, the system can restore a complete KV checkpoint associated with an accepted token prefix and replay tokens accepted after that prefix~\cite{strati2024dejavu,cao2026lumen}. A complete checkpoint contains the historical KV tensors and attention metadata required by all model layers and model-parallel shards. A partial KV copy may accelerate restoration but cannot independently provide recovery coverage.

\section{Motivation and Challenges}
\label{sec:motivation}

The resulting management problem is to decide which KV to retain and
in which tier. For a multi-agent service, the benefit of these decisions
depends on how retained state supports both normal execution and workflow
completion after a failure. Existing fault-tolerant serving and
workflow-aware KV management provide useful mechanisms, but their
objectives leave this joint decision unresolved.

We measure the application-level outcome by \emph{workflow completion
 time} (WCT), $\mathrm{WCT}_w=c_w-a_w$, where $a_w$ is workflow arrival
 and $c_w$ is the time all required outputs become ready. This includes
 model and tool execution, queueing, and dependency waiting.

\subsection{Limitations of Existing Approaches}
\label{sec:motivation-limitations}

\begin{figure*}[t]
  \centering
  \captionsetup[subfigure]{font=normalsize,labelfont=normalfont,
    textfont=normalfont,justification=centering,singlelinecheck=false}
  \begin{subfigure}[t]{2.48in}
    \vspace{0pt}
    \begin{minipage}[c][1.65in][c]{\linewidth}
      \centering
      \includegraphics[width=\linewidth]{figures/workflow.pdf}
    \end{minipage}
    \caption{Four-branch workflow}
    \label{fig:motivation-workflow}
  \end{subfigure}\hfill
  \begin{subfigure}[t]{1.86in}
    \vspace{0pt}
    \begin{minipage}[c][1.65in][c]{\linewidth}
      \centering
      \includegraphics[width=\linewidth]{figures/workflow-gaps.pdf}
    \end{minipage}
    \caption{Workflow-maximum gaps}
    \label{fig:motivation-branch-gap}
  \end{subfigure}\hfill
  \begin{subfigure}[t]{2.52in}
    \vspace{0pt}
    \begin{minipage}[c][1.65in][c]{\linewidth}
      \centering
      \includegraphics[width=\linewidth]{figures/recovery-exposure.pdf}
    \end{minipage}
    \caption{Recovery cost and WCT}
    \label{fig:motivation-recovery-exposure}
  \end{subfigure}
  \caption{Workflow timing and conditional recovery impact.
    (a) Four trace-derived parallel agents A1-A4 with schematic Main agents A0/A5.
    Bars preserve timing; dashed tails assume an all-branch join.
    (b) Empirical CDF of workflow-level gaps over 121 workflows. For each parallel group with temporally inferred fan-in, we subtract the shortest agent execution time from the longest, then retain the largest gap per workflow. 
   (c) Each horizontal bar represents one agent (A1--A4) and shows its
extra delay to reconstruct KV after a failure, measured by full
prefill with Qwen2.5-7B-Instruct-1M. The hatched portion is absorbed
by slack; the solid portion is the calculated WCT increase.}

  \label{fig:motivation}
\end{figure*}

\paragraph{Limitation 1: Fault tolerance without workflow-aware priorities.}
Fault-tolerant serving systems preserve state and coordinate recovery
for interrupted requests. D\'ej\`aVu streams KV checkpoints, and GhostServe
uses host-memory parity to reduce reconstruction
work~\cite{strati2024dejavu,jayakody2026ghostserve}.
LUMEN further coordinates checkpoint placement, recovery routing and
capacity restoration according to worker load~\cite{cao2026lumen}.
These mechanisms address state loss and recovery contention, but do not
use multi-agent dependencies to determine which requests' faster
recovery would advance workflow completion. Request recovery cost and
worker load alone cannot express this benefit.

\label{sec:motivation-workflow}
Figure~\ref{fig:motivation}(\subref{fig:motivation-workflow}) illustrates
why workflow position matters. A2 finishes last, 151.6\,s after the fork,
whereas A1 finishes at 60.4\,s. If the following Main stage requires all
four outputs, A1 can absorb 91.2\,s of additional delay before postponing
the join, whereas any delay on A2 postpones it. This example comes from the Claude Code traces~\cite{semianalysis2026cctraces} and has the fork--join structure illustrated in Figure~\ref{fig:serving_model}. We call the additional
delay a branch can absorb without postponing the required output its
\emph{slack}. 

This situation where agents with parallel execution have different completion times is very common in multi-agent applications, and the time differences can be significant. We analyze the Claude code trace~\cite{semianalysis2026cctraces} related information and visualize this time difference in Figure~\ref{fig:motivation}(\subref{fig:motivation-branch-gap}).
This figure reports each workflow's maximum longest-minus-shortest branch-duration gap,
with a median of 138\,s and a 90th percentile of 1,173\,s. When the Main agent waits for all required branches, differences in branch completion times can leave shorter branches with slack. Additional delays within this slack do not postpone the join or
increase WCT.

These differences in slack affect how KV reconstruction delays workflow
completion after a failure.
Figure~\ref{fig:motivation}(\subref{fig:motivation-recovery-exposure})
shows reconstruction times measured by full prefill with
Qwen2.5-7B-Instruct-1M using sparse
attention~\cite{qwen2025qwen25oneM,jiang2024minference},
and their calculated impact on WCT.
We consider each agent's interruption separately, with other execution
unchanged and the Main agent waiting for all four branches.
A1's 33.2\,s reconstruction delay is fully absorbed by its 91.2\,s slack,
leaving WCT unchanged. In contrast, A2 has no slack, so its 32.6\,s
reconstruction delay increases WCT by the same amount.
Thus, similar reconstruction costs can have very different effects
on workflow completion. Recovery priorities should therefore account
for a request's position in the workflow, in addition to its
reconstruction cost.

\paragraph{Limitation 2: Workflow-aware KV management ignores failures.}
Workflow-aware serving systems already exploit the dependencies
highlighted above, but to accelerate normal execution. One line of work
hands workflow structure to the scheduler: Parrot exposes inter-request
dependencies for application-level scheduling~\cite{lin2024parrot}, Ayo
orchestrates applications as primitive-level dataflow
graphs~\cite{tan2025ayo}, Autellix and SAGA schedule whole agent programs
rather than individual calls~\cite{luo2025autellix,guo2026saga}, and
Cortex and Murakkab pool and configure resources from the same
structure~\cite{pagonas2025cortex,chaudhry2026murakkab}. Another line
manages KV with it directly: InferCept and Continuum decide whether KV is
kept, swapped or discarded during a tool call, and for how
long~\cite{abhyankar2024infercept,li2025continuum}; TokenCake offloads KV
during tool stalls, prefetches it before the next call, and reserves GPU
memory for critical-path agents~\cite{bian2025tokencake}; KVFlow,
CacheScout and PBKV order eviction and prefetching by predicted upcoming
agent invocations~\cite{pan2025kvflow,zhang2026cachescout,zheng2026pbkv};
Pythia predicts future steps and prompt contents to stage KV ahead of
use~\cite{yu2026pythia}; and TOPAS jointly selects resident prefixes and
admitted requests from remaining-path progress, reuse, and movement or
preemption costs~\cite{ni2026topas}. These policies value retained KV by
the requests expected to read it next, and assume that the worker holding
it keeps running.

\begin{figure}[t]
  \centering
  \includegraphics[width=\columnwidth]{figures/gpu-server-availability.pdf}
  \caption{GPU-server availability reconstructed from InfiniteHBD's
    production-training trace~\cite{shou2025infinitehbd}: 400 servers minus
    those within recorded outage intervals. Records include GPU, network,
    power/cooling, and software faults; unknown causes, tests,
    troubleshooting, and changes are also retained. Overlapping intervals
    per server are merged. Availability is measured at the node level.}
  \label{fig:gpu-server-availability}
\end{figure}

Yet failures are not rare, as Figure~\ref{fig:gpu-server-availability}
shows for a production training cluster: 584 recorded unavailability
events affect 231 of 400 servers in about 349 days~\cite{shou2025infinitehbd}. When a server fails, the agent requests running on it are
interrupted together with their GPU-resident KV and any copies swapped
to that host's
memory~\cite{abhyankar2024infercept,pan2025kvflow,bian2025tokencake}.
What remains on other workers or in shared storage are the fixed prompts
and shared prefixes selected by their next expected
read~\cite{pan2025kvflow,yu2026pythia,zhang2026cachescout}; the
interrupted request's history up to the failure is not among them. The
Claude Code trace bounds how much such surviving prefixes could help:
same-agent prefixes cover 97.9\% of all request input (98.7\% for
main-agent and 93.9\% for subagent requests), whereas prefixes from other
agents cover at most 10.1\% (4.4\% and 36.2\%), largely overlapping the
former (Figure~\ref{fig:agent-prefix-reuse}). Most of an interrupted request's
input is thus its own agent's history, which resided with that agent on
the failed worker; it must be reconstructed from the prompt and token
log, at the costs of Figures~\ref{fig:motivation-recompute}
and~\ref{fig:motivation-recovery-exposure}, and which agent it lands on
decides whether that time is absorbed by slack or added to WCT. The
root cause is that failure is absent from these systems' model: whether
a KV object is worth keeping, where, and for how long is decided on the
premise that its worker does not fail. Swapping to host memory frees GPU
capacity rather than escaping a failure, shared storage holds completed
contexts for later reuse rather than a running request's history, and a
copy that would be read only after a failure is worth nothing in their
retention scores. So although these systems already hold the information
needed to tell which agent cannot afford this reconstruction---the
dependency graphs, next-call predictions and remaining paths they extract
are exactly what slack is derived from, and TokenCake even uses them to
reserve GPU memory for critical-path agents~\cite{bian2025tokencake}---that
information is never used to decide, before a failure, for which agent,
at which tier, and for how long a resumable copy should be kept.

\begin{figure}[t]
 \centering
 \includegraphics[width=\columnwidth]{figures/agent-prefix-reuse.pdf}
 \caption{Token-weighted potential prefix coverage in the Claude Code
 trace~\cite{semianalysis2026cctraces}: all 98,827 requests, split into
 56,798 main-agent and 42,029 subagent requests. For each request,
 coverage is the longest input prefix shared with an earlier completed
 request of the same workflow and model, summed over tokens. Same- and
 other-agent bars overlap: other agents add only 0.37\%, 0.01\% and
 2.05\% beyond the best same-agent prefix. These are reuse opportunities,
 not runtime cache hits.}
 \label{fig:agent-prefix-reuse}
 \Description{Three groups show all requests, main-agent requests and
 subagent requests. Same-agent coverage is 97.9, 98.7 and 93.9 percent;
 other-agent coverage is 10.1, 4.4 and 36.2 percent, respectively.}
\end{figure}

\label{sec:motivation-long-context}
The two limitations are mirror images. Fault-tolerant serving decides
\emph{how} to preserve and restore state, but ranks interrupted requests
by recovery cost and worker load. Workflow-aware caching decides
\emph{which} state a workflow will read next, but assumes the worker
holding it survives. Neither answers the question a multi-agent service
faces when a worker fails: which agents' histories must be recoverable
quickly for their workflows to complete on time, and how much of the
finite HBM, DRAM and storage that protection may take from normal reuse.
Combining the two is not a matter of adding slack to a recovery router
or a checkpoint path to a cache. Both objectives draw on the same tiers,
the state each values differs, and both change while agents run.
AegisKV takes this joint decision as its problem. It assumes checkpoint
transport and tiered stores as in the systems above, and contributes the
placement policy---which KV objects each tier holds, on whose behalf,
and when they move---together with the bounded controller and safe
transitions that keep that policy correct online. Realizing it raises
three challenges.

\subsection{Design Challenges}
\label{sec:motivation-challenges}

\paragraph{Challenge 1: Slack requires predicting each agent's remaining execution time online.}
Slack is a statement about the future: how much longer an agent's
branch may run before it, rather than its slowest sibling, gates the
join. The examples in Section~\ref{sec:motivation-limitations} read it
off completed traces. Online, a protection decision needs the remaining
execution time of every running agent at the moment it is made---how
many model calls and tool waits are left, and how long each will
take---and none of this is observed. Tool durations are outside the
serving system's control, the number of further turns depends on outputs
not yet generated, and the branch that gates the join changes as
siblings finish, so one agent's slack shifts whenever another agent
progresses. The quantity to be predicted also spans a wide range:
workflow-maximum gaps run from seconds to over a thousand seconds
(median 138\,s, 90th percentile 1,173\,s in
Figure~\ref{fig:motivation}(\subref{fig:motivation-branch-gap})), so no
fixed rule separates urgent from tolerant branches, and neither recency
nor reconstruction cost stands in for the prediction: A1 and A2
reconstruct in similar time, yet only A2 delays the join. The system
must therefore estimate remaining time from causal features of partial
execution alone, refresh the estimate as agents advance, and derive
priorities that remain useful under prediction error and fall back
safely when the estimate is invalid.

\begin{figure}[t]
  \centering
  \includegraphics[width=\columnwidth]{figures/recovery-cost.pdf}
  \caption{Measured KV reconstruction and reuse costs for
    Qwen2.5-7B-Instruct-1M on two A100 PCIe 40\,GB GPUs
    (TP2, BF16, batch one), with sparse attention above 32K.
    Full prefill starts from empty KV; reuse starts from $(C-1)$ cached
    tokens and computes one. KV reconstruction time includes any
    retrieval and computation until all $C$ tokens' KV is ready.}
  \label{fig:motivation-recompute}
  \Description{Four measured KV reconstruction-time curves for Qwen2.5-7B-Instruct-1M
    from 8K to 1,010,000 tokens: blue circles for resident HBM reuse,
    orange squares for DRAM reuse, purple diamonds for shared-storage
    reuse, and red triangles for full prefill. Both axes are logarithmic.
    At the longest context, the four paths take 1.305, 2.663, 10.156 and
    571.21 seconds, respectively.}
\end{figure}

\paragraph{Challenge 2: KV priority must combine reuse value with recovery benefit, then decide placement.}
Existing systems rank KV by one criterion: workflow-aware caching by the
requests expected to read it next, fault-tolerant serving by the cost of
reconstructing it (Section~\ref{sec:motivation-limitations}). Neither is
the value a multi-agent service needs. A KV object is worth keeping for
two reasons that coexist and can disagree. Normal reuse saves prefill on
the agent's next call whether or not its worker fails. Recovery saves
reconstruction only if a failure occurs, and matters for WCT only to the
extent that the saved time exceeds the agent's slack from
Challenge~1. High slack must therefore not delete history that later
normal calls would recompute, and a low reconstruction cost must not
demote the one agent whose reconstruction would land on the critical
path. Recovery benefits also do not add up: one worker failure
interrupts several co-located agents whose recoveries share readers,
links and target GPUs, and accelerating one of two late branches leaves
their join unchanged. The redefined priority then has to decide storage
management. Retaining a history in any tier removes reconstruction: at
1M tokens, 571.21\,s becomes 10.156\,s from shared storage
(Figure~\ref{fig:motivation-recompute}). A faster tier only shortens
retrieval, to 2.663\,s in DRAM and 1.305\,s in HBM, at a very different
price: that history occupies 57.92\,GB, over 70\% of the two A100s'
combined HBM, and because interrupted agents' histories are largely
unique (Figure~\ref{fig:agent-prefix-reuse}), sharing does not shrink
this cost. A history need not move as a unit, however. Long-context
models increasingly select blocks, tokens or compressed entries per
query (Figure~\ref{fig:sparse-releases}), and their accesses
concentrate: in NOSA-8B, NSA-8B and MiniCPM4.1-8B, about 20\% of
historical blocks receive 87.64\%, 74.76\% and 56.54\% of accesses
(Figure~\ref{fig:design-insights}). A small hot subset in a fast tier
could serve most reads while the complete history stays in cheaper
tiers, provided the subset is chosen from causal history alone and a
missed block remains retrievable. Placement must therefore turn the
combined priority into decisions per tier and per block---which objects
each tier holds, on whose behalf and for how long---separating
retention that avoids reconstruction from acceleration that shortens
retrieval, within finite tier budgets.

\begin{figure}[t]
 \centering
 \includegraphics[width=\columnwidth]{figures/sparse-models.pdf}
 \caption{Selected sparse/hybrid model releases and reported context
 lengths. IDs: 1 MiniCPM4~\cite{minicpm4model};
 2 MiniCPM4.1~\cite{openbmb2025minicpm41};
 3 DeepSeek-V3.2-Exp~\cite{deepseek2025v32exp};
 4 DeepSeek-V3.2~\cite{deepseek2025v32model};
 5 GLM-5~\cite{zai2026glm5model};
 6 MiniCPM-SALA~\cite{openbmb2026sala};
 7 DeepSeek-V4 preview~\cite{deepseek2026v4model};
 8 MiniMax-M3~\cite{minimax2026m3};
 9 GLM-5.2~\cite{zai2026glm52};
 10 LongCat-Flash-Lite-Sparse~\cite{longcat2026sparse}.
 Dates are announcements, or first official repository dates for
 IDs 5, 9 and 10. Markers identify the sparse branch's selection unit,
 including SALA's block branch and V4's compressed sparse branch;
 they do not imply uniform attention across layers. This is a selected
 survey, not an adoption-share or sparsity-rate estimate.}
 \label{fig:sparse-releases}
 \Description{Ten selected sparse or hybrid attention releases between
 June 2025 and July 2026, with reported contexts from 32K to 1M.
 Circles, squares and diamonds distinguish block, token and compressed
 entry selection. This is a literature survey, not a sparsity-rate curve.}
\end{figure}

\begin{figure}[t]
 \centering
 \includegraphics[width=\columnwidth]{figures/kv-access-concentration.pdf}
 \caption{Native KV access concentration for NOSA-8B~\cite{huang2026nosav2},
 NSA-8B (MiniCPM-4-8B-NSA)~\cite{fasa2026minicpm4nsa}, and
 MiniCPM4.1-8B~\cite{openbmb2025minicpm41}. We use the final 512 queries
 of matched 32,704-token prefixes from three PG19 texts. Each query
 counts a historical 64-token block once across heads. Rank the 503
 blocks within each layer/window; average access shares over all
 32 layers, then texts. Bands show text min--max; the diagonal denotes
 uniform access. Compressed/selector state is excluded. This is
 retrospective concentration, not prediction accuracy or recovery speedup.}
 \label{fig:design-insights}
 \Description{Cumulative historical block access curves for NOSA-8B,
 NSA-8B and MiniCPM4.1-8B, averaged over every layer and three texts.
 About 20 percent of blocks account for 87.64, 74.76 and 56.54 percent
 of accesses, respectively; all three curves exceed the uniform baseline.}
\end{figure}

\paragraph{Challenge 3: Extra protection must not disrupt normal inference.}
Every copy kept for recovery beyond what normal serving would keep is
extra work on the serving path. Writing checkpoints and moving them
between tiers consumes the DRAM and storage bandwidth that normal reuse
and concurrent recoveries read from, and a copy held in HBM occupies
memory that inference batches need: every HBM byte spent on protection
is taken from serving. Deciding what to protect is not free either.
Every generated token extends the unsaved tail, every tool return and
new request shifts slack, and every batch changes the HBM inference
needs, so re-scoring all resident objects on each event is prohibitive
at block granularity, while stale scores keep the wrong state in the
fast tiers. Protection must also not undermine itself: transfers are
asynchronous and the failure they guard against can strike midway;
releasing a DRAM copy before its storage replacement commits, or
counting an object as local after its owner failed, silently leaves a
checkpoint without a recoverable copy. The controller therefore needs
bounded work per event, a protection capacity that yields to inference
when foreground demand rises and grows back when it falls, migration
paced within the bandwidth budget, and transitions that keep every
committed checkpoint with a reachable owner.

\label{sec:motivation-goals}
These challenges organize the design. Workflow Tracker predicts remaining
execution time and derives slack from it for Challenge~1
(Section~\ref{sec:workflow-protection}); the slack-aware request
priority, block-level demand and renewable per-agent credit define the
combined priority and its placement for Challenge~2
(Sections~\ref{sec:placement-value}--\ref{sec:sparse-placement}); and
the lazy event loop, the AIMD Adapter for Peer capacity and versioned
transfers keep protection off the inference path for Challenge~3
(Sections~\ref{sec:placement-plan} and~\ref{sec:adapter}).

\section{AegisKV Design}
\label{sec:design}

AegisKV allocates finite KV storage to reduce workflow completion time
(WCT) while preserving normal cache reuse. This section presents the
architecture, then each component in the order of the challenges in
Section~\ref{sec:motivation-challenges}.

\subsection{System Overview}
\label{sec:system-overview}

AegisKV combines workflow priorities with a three-tier KV hierarchy
(Figure~\ref{fig:system-overview}). Peer HBM keeps selected
layers and hot blocks on another data-parallel (DP) replica for fast
recovery; shared DRAM and storage retain the remaining state for reuse
and recovery. Together, the three tiers cover each committed checkpoint,
with one managed owner per object after transfers finish. This division lets
the system retain long histories while spending scarce HBM on the state
of zero-slack requests with bounded per-request quotas.

The example uses two DP replicas, each with two tensor-parallel (TP)
ranks. In this compatible head-sharded layout, the ranks hold different
KV heads for the same token positions. Peer placement maps selected
objects of source shard $i$ to rank $i$ on the takeover DP.

The Application supplies requests and workflow progress to the Controller.
\textbf{(1)} Workflow Tracker estimates how much delay each Agent can
absorb, while \textbf{(2)} Recovery-aware Router selects its execution and
takeover DPs. Using this urgency and location information, \textbf{(3)} KV
Master decides where to retain each KV object and which hot blocks merit
Peer space. Within the Engine, \textbf{(4)} Adapter adjusts the available
Peer capacity as serving demand changes. Connector executes Master's
transfer plans and reports their outcomes. If a worker fails, Router redirects its requests to the
takeover DP, where Connector combines Peer hits with lower-tier reads to
resume execution.

\begin{figure}[t]
 \centering
 \includegraphics[width=\columnwidth]{figures/system-overview.pdf}
 \caption{AegisKV with DP=2 and TP=2; callouts (1)--(4) match the text.
 Each TP row is a GPU/HBM rank. For DP0-to-DP1 recovery, DP1's rank $i$
 keeps selected layer/block objects of DP0's shard $i$; the reverse
 direction is symmetric. Local KV is execution state; Peer holds managed
 recovery objects. Shared DRAM/storage hold the remainder outside both
 DP boundaries. Blue links carry KV, dark dashes control, red dashes
 retry, and orange blocks unsealed tails. The Master--Connector link
 shows plans and returned status on DP1/TP0; Adapter is grouped per DP.}
 \label{fig:system-overview}
 \Description{Three layers with numbered Tracker, Router, KV Master and
 Adapter. Two side-by-side DP replicas each contain TP0 and TP1, whose
 local and Peer HBM hold distinct KV head shards 0 and 1. Rank-level
 Connectors exchange corresponding shards across DPs and access shared
 DRAM and storage. A separate Master control link carries plans and status.}
\end{figure}

\subsection{Workflow Tracker: Remaining Time and Slack}
\label{sec:workflow-protection}

To make workflow priorities available online, Tracker estimates remaining
execution before deriving slack. It maintains a residual DAG whose nodes
are Agent segments between exposed fork/join boundaries or completion.
A segment can span several LLM calls and tool waits. Child execution
appears on child paths, not again in the parent; unknown future forks
are excluded.

A GBDT predicts segment seconds remaining, $d$, from causal features: role/model,
current input length, elapsed time, completed-call count and previous
calls' lengths, timings and context growth. Missing history is masked;
training, calibration and evaluation split whole workflows. Labels end
at the segment boundary, and future call outcomes are never features.

Algorithm~\ref{alg:workflow-tracker} takes residual DAG $G$ and features
$X$. Forward/backward critical-path passes~\cite{kelley1959cpm} compute
earliest finishes $EF$ and latest finishes $LF$ relative to the predicted
workflow finish. Their difference is nominal slack. The guard $g$ is an
upper empirical quantile of positive slack overestimation on calibration
workflows; subtracting it gives advisory slack $s$, not a per-request
guarantee. Requests in one segment share its margin.

\par\medskip\noindent\begin{minipage}{\columnwidth}
\raggedright
 \refstepcounter{designalgorithm}\label{alg:workflow-tracker}
 \textbf{Algorithm \thedesignalgorithm: Remaining time and slack}
 \par\smallskip
 \begin{tabularx}{\columnwidth}{@{}r@{\quad}X@{}}
 \toprule
 1 & $(G,X)\gets\textsc{SnapshotProgress}()$ \\
 2 & $\nu\gets\textsc{NextVersion}()$ \\
 3 & \textbf{if not} \textsc{Supported}$(G,X)$ \textbf{then} \\
 4 & \quad \textsc{PublishInvalid}$(\nu)$; \textbf{return} \\
 5 & $d\gets\max(0,\textsc{GBDT}(X));\quad g\gets\textsc{Guard}()$ \\
 6 & $EF\gets\textsc{Forward}(G,d)$ \\
 7 & $LF\gets\textsc{Backward}(G,d,\max_v EF_v)$ \\
 8 & \textbf{for} $v\in V(G)$ \textbf{do} \\
 9 & \quad $s_v\gets\max(0,LF_v-EF_v-g_v)$ \\
 10 & \textsc{Publish}$(\nu,G,d,g,s)$ \\
 \bottomrule
 \end{tabularx}
\end{minipage}\par\medskip

\textsc{Supported} checks inputs, dependencies, calibration and freshness;
invalid estimates grant no positive slack credit. Progress events and
periodic refresh publish versioned durations $d$, guards and slack to
Master, invalidating older priority estimates. The DAG passes cost
$O(|V|+|E|)$. Section~\ref{sec:eval-prediction} reports available prediction
diagnostics and specifies residual-segment validation.

\subsection{Recovery-aware Router}
\label{sec:recovery-router}

Router establishes where normal and resumed execution run. Normal
dispatch balances load and compatible prefix locality.
Figure~\ref{fig:agent-prefix-reuse} reports 98.7\% and 93.9\% potential
same-agent token coverage for main-agent and subagent requests,
respectively. Agent affinity is a preference: prefixes from other Agents
also offer reuse, especially for a subagent's first request.

At arrival, Router records a versioned takeover binding to a DP with a
compatible layout outside the source failure domain. Its corresponding
tensor-parallel (TP) shards can own selected Peer objects. On confirmed
failure, Router fences the old epoch and activates that target subject
to workspace and admission constraints. A new target cannot count the
old target's Peer objects as local: rebinding retrieves reachable objects
or uses an older checkpoint/replay fallback if their sole owner was lost.
This binding gives Master concrete access paths and target load for
recovery profiling and Peer placement.

\subsection{KV Master: Workflow-aware Tiered Placement}
\label{sec:kv-master}

Master first combines request urgency with within-request KV demand,
then carries this priority across repeated and shared uses. Renewable
credit manages lower-tier retention; a fixed-size, request-level policy
allocates Peer acceleration. An event loop connects both policies to
actual transfers and available capacity.

\subsubsection{Slack- and Access-aware KV Priority}
\label{sec:placement-value}
\label{sec:prefix-value}

Two factors determine a request's interest in a KV block:
how strongly recovery affects its workflow, and how likely execution is
to need that block. Let $R_r$ be the profiled recovery delay under the
request's current placement, including lower-tier reads, replay and
queues, and $s_r$ its guarded slack. We use the exposed fraction
\begin{equation}
 u_r=\frac{[R_r-s_r]^+}{\max(R_r,\varepsilon)},
 \qquad P_{r,l,b}=u_r\,\widehat p_{r,l,b},
 \label{eq:value}
\end{equation}
where $\varepsilon>0$ avoids division by zero. Thus zero slack gives full
urgency when recovery takes time; slack covering recovery removes this
urgency, not ordinary reuse value. Invalid Tracker estimates use zero
slack. $P$ is a ranking signal, not an additive estimate of WCT saved.

Within one dense-attention layer, every causally visible historical block
is needed by subsequent queries, so its access factor is one. Earlier
creation alone does not imply greater future value. Sparse attention
instead uses a recent window of completed queries:
\begin{equation}
 \widehat p_{r,l,b}=
 \frac{\text{queries selecting block }b}
      {\text{queries for which block }b\text{ was visible}}.
 \label{eq:selection-probability}
\end{equation}
Each query counts a physical block once across heads. Visibility is
defined before sparse selection, not after candidate pruning. A block
selected 40 times in 200 opportunities has lower estimated demand than
one selected 15 times in 30; raw counts reverse that ranking. A calibrated
layer prior smooths small denominators. This causal estimate predicts
demand without modifying the native selector; mandatory auxiliary state
remains a dependency regardless of its measured frequency. The access
concentration in Figure~\ref{fig:design-insights} is what makes per-block
demand worth estimating; the estimate stays causal, and a block missing
from a fast tier is fetched from its owner.

Priority and retention cost are distinct. For equal-size, equal-cost
blocks in the same dense layer, the access contribution is equal; object
sizes, transfer paths and reconstruction costs can still differ.
Figure~\ref{fig:retention-value} illustrates why cost per byte also matters.
This motivates cost-aware lower-tier credit, while the Peer policy below
deliberately uses a simpler urgency rule.

\begin{figure}[t]
 \centering
 \includegraphics[width=\columnwidth]{figures/retention-value.pdf}
 \caption{Slack alone does not rank retention value. In this isolated-fault
 example, A uses two capacity units to reduce recovery from 4 to 2\,s;
 B uses one to reduce 18 to 8\,s. A at slack 0 saves 1\,s/unit, whereas
 B at slack 6 saves 10\,s/unit. Normal-use value is zero. This constructed
 example motivates cost-aware lower-tier retention and exposes a limit
 of urgency-only Peer admission; it is not a policy-speedup result.}
 \label{fig:retention-value}
 \Description{A higher-slack object can save more workflow time per byte.}
\end{figure}

\subsubsection{Renewable Credit with Lazy Updates}
\label{sec:renewable-credit}

The preceding priority concerns one request. Across requests, a shared
prefix can serve several Agents, while one Agent can return after a tool
wait. Master gives each interested Agent its own share of an object's
credit. A different Agent adds a share; the same Agent's new use refreshes
its share instead of accumulating another copy of its value. Agent
identity includes its workflow instance, and physical capacity is charged
only once.

For a lower-tier object $x$, the current grant combines a normal-reuse
term with its recovery priority:
\begin{equation}
 g_{a,x,j}=v^{\mathrm{reuse}}_{a,x,j}
                +\lambda_f c_{a,x,j}P_{a,x}.
 \label{eq:credit-grant}
\end{equation}
Here $j$ is DRAM or storage; $\lambda_f$ weights fault risk over a common
decision horizon. The profiled cost $c$ measures the additional service
work if the object leaves $j$: storage access after DRAM demotion, or
reconstruction after final storage removal. The reuse term multiplies a causal next-use estimate by the corresponding
saved service work, independently
of slack. This preserves useful state for long-slack, multi-round Agents.
When recovery requires a dependent group, its profiled cost is distributed
by bytes rather than repeatedly assigned to every block. Summed Agent
shares express competing interests, not exact joint WCT: simultaneous
delays at a workflow join are not additive.

A one-time grant becomes stale; granting fresh credit at every checkpoint
can retain unused state indefinitely. Inspired by Landlord and
GreedyDual-Size~\cite{young2002landlord,cao1997greedydual}, Master separates
current grant $g$ from already-paid rent $d$:
\begin{equation}
 C_{x,j}=\sum_a[g_{a,x,j}-d_{a,x,j}]^+,
 \qquad H_{x,j}=\rho_j+C_{x,j}/b_x.
 \label{eq:rental-credit}
\end{equation}
The pool pressure clock $\rho_j$ charges bytes $b_x$, rather than elapsed
time. Increasing it by $\delta$ removes $\min(C_{x,j},b_x\delta)$, split
proportionally among positive Agent shares. All committed residents age,
including leased objects; eviction eligibility is separate. Actual reuse or recovery consumption renews only the using Agent,
once per object and request/use ID. Attention observations refresh demand;
checkpoint growth, score updates and migration preserve paid rent.
Expired Agent interest is removed. Figure~\ref{fig:renewable-credit}
illustrates these different events. Recency enters through renewal,
demand through selection probability and reuse history, and workflow
urgency through $u_r$; these are explicit signals, not a claim to subsume
every LRU/LFU variant.

\paragraph{Avoiding a credit update per token.}
The engine advances one query counter per layer and records only selected
block IDs. A block's birth index and the window boundary determine its
opportunities without incrementing every block's denominator. Bounded
query buckets expire old selected counts; dense attention needs no such
counters. Statistics are batched at checkpoint events. Slack changes
update a request scalar and version, leaving its blocks untouched.

Master materializes a block only on use, candidate inspection or a
rotating scan. It first settles rent since that block's last clock,
then reads current request/demand versions and refreshes the grant while
retaining paid amounts. Grants remain fixed between materializations;
the shared pressure clock therefore avoids scanning all residents on
each charge. Refreshing only a stale heap minimum is insufficient when
unseen keys can decrease. Each round refreshes $q$ low-key candidates
and $q$ rotating alternatives, within a fixed work budget, before choosing
victims. This is approximate selection, compared with eager refresh in
Section~\ref{sec:eval-rental}; it is not an exact global minimum. Work
scales with observed selections, examined candidates and their Agent
shares, rather than a full-cache update per generated token.

\subsubsection{Greedy Peer Protection}
\label{sec:sparse-placement}

Peer accelerates recovery without participating in rent. Because accesses
concentrate on a subset of blocks (Figure~\ref{fig:design-insights}), a
fixed hot set can serve most reads of the selected layers. For each supported
model/layout, calibration fixes checkpoint interval $Z$, prefix depth $L$
and $K$ blocks per selected layer. The online policy changes hot-block
identities and request membership, not these three parameters. A Peer
slot reserves at most $LK$ blocks plus a separately budgeted mandatory
state allowance. Short histories use fewer blocks; shared physical
objects are charged once. This is a fixed quota, not a claim of identical
physical occupancy for every request. Engine block size is configured
separately from $Z$; a measurement's 64-token unit is not universal.

Figure~\ref{fig:sparse-recovery} shows the resulting placement. Each sealed
layer/block/TP object has one managed owner, distinct from execution
working KV. Peer, DRAM and storage together cover a committed checkpoint;
moving an object to Peer releases its lower-tier copy after handoff.
Unselected blocks stay below and demand misses fetch their actual owners.

\begin{figure}[t]
 \centering
 \includegraphics[width=\columnwidth]{figures/sparse-recovery.pdf}
 \caption{A checkpoint across three exclusive tiers, with a fixed Peer
 quota $L=2$, $K=2$. Each cell groups a layer/block's TP objects; four
 cells reside in Peer and the others in DRAM/storage. Layer 1 needs
 B0/B1/B2: B0 is local, B1 needs H2D, and B2 needs a storage read plus
 H2D. The dashed tail is uncommitted. Hot identities can change within
 the quota; a Peer victim must migrate down before its space is reusable.}
 \label{fig:sparse-recovery}
 \Description{Three-layer ownership map with two Peer blocks in each of
 the first two layers and actual-owner recovery paths.}
\end{figure}

Only requests with clipped guarded slack $s_r=0$ enter the Peer candidate
queue, including requests whose slack reaches zero while running.
Admitted requests select $K$ blocks per layer by $\widehat p$, with recency
then block ID breaking ties. Completion or tool handoff makes that
request's slot evictable. On admission or quota shrink, Master first
reclaims ended requests, then active requests in decreasing slack order.
Ordinary admission preserves an incumbent at equal slack; a new tied
request waits for Peer while retaining lower-tier protection. Required
shrink breaks active ties by least recent request use. Active requests
therefore have a policy preference, not an absolute eviction lock;
physical read/copy leases are always respected. A shared object is released
only when no retained slot needs it. Different takeover DPs cannot both
count its sole owner as local.

\paragraph{Why fix a partial prefix of layers?}
Let $M$ be committed history and $N$ its acknowledged, uncommitted tail.
Recovery replays $N$; that computation can overlap later-layer reads.
Peer removes only exposed transfer waiting, beyond the reconstruction
already avoided by lower-tier state. Figure~\ref{fig:design-pipeline}
includes storage reads, H2D and replay on one clock. Demand misses make
sparse loading slower than full-layer pipelining in this example; Peer
removes about 15\,ms relative to the same sparse demand. The comparison
illustrates partial residency, not a universally optimal $L$.
Calibration sweeps fixed quotas over declared history/tail and load
profiles, selects a configuration under an HBM budget, then freezes it
before evaluation on other workflows. Heterogeneous requests may prefer
different depths; planned adaptive-depth and small-instance oracle
comparisons will quantify the cost of this simplification.

\begin{figure}[t]
 \centering
 \includegraphics[width=\columnwidth]{figures/recovery-pipeline.pdf}
 \caption{Analytical storage/H2D/replay schedules: 32 layers, 32K history,
 64-token tail, and 1,024\,MiB unique KV. (a)--(c) split history equally
 between DRAM/storage; (d) moves $L=7$, $K=336$ 64-token blocks per layer
 to Peer: 147\,MiB Peer, 438.5\,MiB per lower tier. The existing
 MiniCPM4-8B replay calibration is 2.870\,ms/layer; storage/H2D rates
 of 3.125/6.25\,GB/s are assumed. Both sparse cases preload 336 blocks
 and miss on 64 per layer; misses become known after the previous layer
 finishes. These demands are constructed. Cold KV follows the plotted
 endpoint. Placement is complete; I/O stages are independent, workspace
 sufficient and foreground traffic absent. These are mechanism schedules,
 not integrated measurements or an optimal quota search.}
 \label{fig:design-pipeline}
 \Description{Four analytical recovery schedules separating storage
 reads, H2D, computation and exposed wait at equal unique coverage.}
\end{figure}

\subsubsection{Online Three-Tier Management}
\label{sec:placement-plan}

Prefill offers sealed blocks; decode reviews every fixed $Z$ tokens, and
completion/tool handoff offers a final snapshot, including short rounds.
Actual use, Slack changes, capacity updates and transfer outcomes also
invoke Algorithm~\ref{alg:placement}. New blocks receive the creating
Agent's initial grant; unchanged prefix objects receive no new grant.

\paragraph{Admission and pressure.}
Free lower-tier space admits new state, including zero-credit state.
Under pressure, Master refreshes candidate victims and trials enough
byte-proportional rent to exhaust an examined, safely removable bundle.
Incoming objects pay the same trial rent; those reaching zero are deferred
from that tier. DRAM rejection can send state directly to storage.
Committed zero-credit objects remain available until space is needed;
DRAM victims migrate to storage, whose separate ledger values retention
against reconstruction. Last-copy removal must publish a valid older
checkpoint or input/log replay fallback for every affected manifest.
Dependencies and physical leases constrain all victim choices.

Peer admission/shrink follows the request order above. A hot-set change
swaps identities within the same slot. Both can require Peer-to-DRAM and
DRAM-to-storage cascades. If a safe destination cannot be reserved, the
move waits and the source remains occupied. Paid amounts persist per
lower pool across round trips, with no rent during absence. Neither a
move through Peer nor a value-only refresh renews credit.

\par\medskip\noindent\begin{minipage}{\columnwidth}
\raggedright
 \refstepcounter{designalgorithm}\label{alg:placement}
 \textbf{Algorithm \thedesignalgorithm: Online three-tier management}
 \par\smallskip
 \begin{tabularx}{\columnwidth}{@{}r@{\quad}X@{}}
 \toprule
 1 & \textsc{ReconcileTransfers}$(e)$; \textsc{PublishSignals}$(e)$ \\
 2 & \textsc{RecordNewBlocksAndUses}$(e)$ \\
 3 & $Q\gets\textsc{LowKeysAndRotation}(q)$ \\
 4 & \textbf{for} $x\in Q$ \textbf{while} \textsc{WorkLeft}() \textbf{do} \\
 5 & \quad \textsc{SettleRent}$(x)$; \textsc{RefreshGrantAndKey}$(x)$ \\
 6 & $P_D\gets\textsc{LowerTierAdmission}(e,Q)$ \\
 7 & $P_H\gets\textsc{PeerSlots}(s,B,L,K,e)$ \\
 8 & $P\gets\textsc{PrepareCascade}(P_D,P_H)$ \\
 9 & \textbf{if} $P=\varnothing$ \textbf{then} \textsc{Defer}$(e)$; \textbf{return} \\
 10 & \textbf{if not} \textsc{ReserveIfCurrent}$(P)$ \textbf{then} \\
 11 & \quad \textsc{Requeue}$(e)$; \textbf{return} \\
 12 & \textsc{Submit}$(P)$ \quad\textit{completion/failure emits a new event} \\
 \bottomrule
 \end{tabularx}
\end{minipage}\par\medskip

Connector copies, confirms the destination, redirects ownership, and
releases the source after readers drain. Both copies count during the
handoff. Rejected previews change neither ownership nor committed rent;
successful capacity handoffs commit their trial rent and residual grants.
Only complete dependency coverage advances $M$; copy lag can make $N>Z$.
The protected source failure must leave advertised owners reachable;
additional owner loss requires an explicit fallback. Capacity checks
validate feasible execution, not global WCT optimality. Dynamic values,
dependent objects and budgeted refresh also prevent inheriting Landlord's
competitive guarantee. Section~\ref{sec:eval-rental} specifies same-value
comparisons with direct revaluation and separate renewal, aging, and
shared-Agent controls, and reports the historical rental diagnostic.

\begin{figure}[t]
 \centering
 \includegraphics[width=\columnwidth]{figures/renewable-credit.pdf}
 \caption{Credit for one shared, unit-size DRAM object. Grants 3 and 2
 become remaining shares 1.8 and 1.2 after rent of 2. Updating A's
 grant to 4 preserves its paid amount 1.2, leaving 2.8; an actual new
 use by A renews its share to 4. B stays at 1.2. Lazy materialization
 settles earlier rent before updating the grant. This is a constructed
 accounting example; Peer has no rental clock.}
 \label{fig:renewable-credit}
 \Description{A per-Agent ledger separating aging, value refresh and renewal.}
\end{figure}

\subsection{Adapter: Guarded AIMD for Peer Capacity}
\label{sec:adapter}

Master selects content within a quota; Adapter adjusts that quota as
foreground inference demand changes. It uses pressure-driven multiplicative
decrease and stable-window additive increase~\cite{chiu1989aimd}.
The total Peer authorization has a physical ceiling
\begin{equation}
 C=\max\{0,\min(B_{\mathrm{cap}},
 B_{\mathrm{pool}}-B_{\mathrm{serve}}-B_{\mathrm{reserve}})\}.
 \label{eq:peer-ceiling}
\end{equation}
The pool excludes weights/fixed allocations; serving high watermark
excludes Peer, and reserve covers growth/workspace. Existing Peer bytes
are not subtracted twice. $B_{\rm cap}$ can be a configured pool fraction.

The current quota $B$ and stable-window count $z$ start at zero.
Pressure (allocation stalls, reserve violations or
memory-related admission waits) multiplies $B$ by $0<\gamma<1$. After $m$
stable windows below a lower wait threshold, deferred zero-slack Peer
demand, counted before quota filtering, permits an additive byte step
$\Delta$. Intermediate windows reset
$z$; missing telemetry disables growth. The ceiling clips every update;
emergencies may decrease within a window, but growth requires completed
stable windows. Parameters and window length are configurable.

\par\medskip\noindent\begin{minipage}{\columnwidth}
\raggedright
 \refstepcounter{designalgorithm}\label{alg:adapter}
 \textbf{Algorithm \thedesignalgorithm: Windowed Peer quota}
 \par\smallskip
 \begin{tabularx}{\columnwidth}{@{}r@{\quad}X@{}}
 \toprule
 1 & $C\gets\textsc{SafeCeiling}();\quad B'\gets B$ \\
 2 & \textbf{if} \textsc{Pressure}() \textbf{then} $(B',z)\gets(\gamma B,0)$ \\
 3 & \textbf{else if} \textsc{Stable}() \textbf{then} \\
 4 & \quad $z\gets z+1$ \\
 5 & \quad \textbf{if} $z\ge m$ \textbf{then} \\
 6 & \qquad \textbf{if} \textsc{PeerDemand}() \textbf{then} $B'\gets B+\Delta$ \\
 7 & \qquad $z\gets0$ \\
 8 & \textbf{else} $z\gets0$ \\
 9 & $B\gets\textsc{RoundDownPage}(\min(C,B'))$ \\
 10 & \textsc{PublishBudget}$(B)$; \textsc{EmitCapacityEvent}() \\
 \bottomrule
 \end{tabularx}
\end{minipage}\par\medskip

Adapter sets the total budget; Master admits fixed-$L/K$ request slots
and refreshes their identities at fixed $Z$ and other placement events. If occupancy plus reservations exceeds a reduced
quota, Master pauses Peer admission and selects objects to demote. Their
DRAM destinations may first require storage demotion. Serving can reuse
the memory only after these transfers and ongoing readers finish; the
new quota is an authorization target, not an instantaneous free-space claim.


\section{Evaluation}
\label{sec:evaluation}

The central question is whether workflow-aware protection reduces
\emph{workflow completion time} (WCT) after a failure without paying a
larger cost during normal service. We organize the evaluation around
end-to-end outcomes, recovery-path costs, mechanism contributions,
prediction quality, and operating limits. The experimental design follows
three useful precedents: application-level outcomes in workflow-aware
serving~\cite{lin2024parrot,ni2026topas}; separate interrupted and
uninterrupted cohorts in failure recovery~\cite{cao2026lumen}; and
end-to-end measurements followed by restoration, scheduling, and
resource-pressure ablations in KV-cache systems~\cite{gao2025hcache,qin2025mooncake}.

\paragraph{What is measured.}
The current artifact contains physical KV-path microbenchmarks,
held-out predictions on observed agent traces, native sparse-access
captures, and historical analytical diagnostics. It does not yet contain
a fault-injection campaign for the current fixed-$L/K/Z$, exclusive
three-tier controller. We report the available evidence at its actual
scope and specify the missing experiments below. In particular, an older
Peer-disabled allocation model and an isolated rental model are
diagnostics, not implementations of the present policy. The figures
distinguish physical measurements, trace observations, and models; no
synthetic curve stands in for an unexecuted system comparison.

\subsection{Methodology and Comparison Contract}
\label{sec:eval-method}

\paragraph{Deployment and workloads.}
The primary system experiment requires at least two compatible serving
replicas in distinct failure domains, independently reachable backing
storage, and a durable input/output log. GPU-process, replica, host, and
backing-store failures are separate cells. Record DP and TP independently,
the surviving compute capacity, physical interconnect, backend, model
revision, precision, attention kernels, block size, and controller version.
Local NVMe, remote NVMe, and a shared filesystem are distinct backends;
a local disk is not host-failure-safe. These deployment fields remain
unfilled until an executable prototype and hardware allocation are frozen.

Use executable coding, multi-hop question-answering, and controlled
sequential/fork--join workflows. Preserve the original DAG once, complete
prompt histories, workflow IDs, and absolute arrival sequence. Frozen
tool responses provide repeatable primary comparisons; live tools are
a separate validity experiment. Dense and native-sparse model paths are
separate cells, with unsupported context lengths explicitly excluded.
Calibrate one fixed baseline's sustainable healthy load, then apply the
same open-loop arrivals to every policy. The initial decisive slice uses
three workload families, two compatible model paths, one load/budget
setting, four policies, two fault conditions, and ten paired campaigns:
480 planned run slots, not 480 completed runs. Broader load, context,
capacity, and fault-position scans follow the decisive slice.

\paragraph{Baselines and resource accounting.}
Recompute retains accepted tokens and application state but no additional
recoverable KV; it measures the benefit of checkpointing itself.
Hierarchical LRU (H-LRU) and reuse-only protection use the same transport,
checkpoint trigger, recovery scheduler, ownership protocol, and total
physical resources as AegisKV. H-LRU tests whether the new priorities beat
a tuned hierarchy; reuse-only isolates recovery urgency from normal
locality. A zero-slack control isolates workflow urgency while preserving
normal-reuse value. Peer-off and fixed-total-Peer-cap controls isolate
Peer acceleration and the Adapter.
All-DRAM and small-instance offline policies are explicitly labeled
references with their relaxed assumptions. Compatible external
D\'ej\`aVu, GhostServe, LUMEN, or Mooncake artifacts provide additional
system comparisons~\cite{strati2024dejavu,jayakody2026ghostserve,cao2026lumen,qin2025mooncake};
a reimplemented mechanism is named as such, and an incompatible artifact
is not assigned invented performance.

All equal-resource policies pay for execution and protection HBM,
metadata, staging, reserved and retiring bytes, checkpoint writes, and
migration traffic. Shared objects consume one stable physical allocation;
both copies count during handoff. The union of Peer, DRAM, and storage
owners covers a checkpoint. We do not give Peer a free complete lower-tier
replica. $L$, $K$, and $Z$ are selected on calibration workflows and frozen
for held-out evaluation. Adaptive per-request depth or hot-set size is
a separate reference; the current Adapter already varies total Peer
capacity while keeping the per-request quotas fixed.

\paragraph{Outcomes and statistical unit.}
For policy $p$ and workflow $w$, use that policy's own fault-free mate:
\begin{equation}
 \begin{aligned}
 D_{w,p}&=\mathrm{WCT}^{F1}_{w,p}-\mathrm{WCT}^{F0}_{w,p},\\
 G_{w,p}&=D_{w,\mathrm{H\mbox{-}LRU}}-D_{w,p}.
 \end{aligned}
 \label{eq:eval-penalty}
\end{equation}
$F0$ denotes no injected failure; $F1$ denotes the matched fault condition.
Positive $G$ indicates a smaller failure penalty. Report absolute WCT,
completion by a common horizon, SLO goodput, and healthy-service cost as
well as $D$ and $G$. Separate directly interrupted workflows, workflows
sharing recovery resources, and an isolated control cohort. Campaigns
are the independent repetition unit; workflows, requests, and blocks
within a campaign are not independent runs. Report both object counts
and campaign counts, and aggregate paired differences within each
campaign before computing campaign-level uncertainty.
TTFT starts at the original request arrival, including retries. For
requests already decoding at failure, report failure-to-first-resumed-token
latency and the inter-token latency tail. End-to-end TPOT includes recovery
pauses between the first and last accepted output tokens; decoder-only
TPOT is a separately labeled diagnostic. Replayed or discarded tokens do
not enlarge the accepted-output denominator.

\subsection{End-to-End Workflow Completion}
\label{sec:eval-workflow}

\paragraph{Decisive experiment.}
The primary fault is triggered when a fixed workflow/Agent/request reaches
80\% of its frozen logical token progress. Kill the complete source
replica and include every interrupted request. Sweep 20/50/80/95\%
progress separately; do not substitute one wall-clock timestamp for
equal progress. In attribution runs, freeze the pre-fault source mapping
and committed coverage. In full-system runs, allow routing and achieved
frontiers to differ, but compare the original arrival cohort and report
the actual interrupted IDs and wall-clock trigger times. Do not select
a different interrupted subset for each policy. Routing and trigger-time
differences remain part of the full-system effect; mechanism attribution
uses the frozen mapping and coverage control above.

The primary result should pair a WCT distribution with same-ID $G$
values for each workload family and load. Show completion and timeouts
alongside latency; use a shared eligible cohort for complete-case
latency. Failure penalties can be zero or negative, so ratios with $D$
as the denominator are inappropriate. A convincing result must improve
over H-LRU or reuse-only, not only over Recompute, and must survive
accounting for the healthy-service cost.

\paragraph{Available allocation diagnostic.}
Figure~\ref{fig:eval-placement-pairs} reanalyzes the existing
36-configuration, four-branch model. It varies state size, branch slack,
and storage bandwidth while comparing identical configurations. The
historical workflow-value allocator improves 17 cases and ties 19. In
the base case, LRU and workflow value add 3.6208 and 1.6125\,s to a
12\,s residual join, respectively. The 2.0083\,s difference is a
55.5\% reduction in the \emph{failure penalty}, not in full-workflow WCT.
This model has no Peer tier, fixed-slot admission, or lazy lower-tier
controller. It supports the relevance of workflow priorities and the
choice of a paired metric, but not an end-to-end AegisKV speedup.

\begin{figure}[tb]
 \centering
 \includegraphics[width=\columnwidth]{figures/evaluation/eval-placement-pairs.pdf}
 \caption{Historical allocation model: paired reduction in failure penalty
 relative to LRU across 36 equally weighted configurations. Positive values
 favor workflow-value placement; 19 ties remain visible. These are
 deterministic Peer-disabled scenarios, not measured workflows.}
 \label{fig:eval-placement-pairs}
 \Description{Paired LRU minus workflow-value failure penalties over all
 36 historical analytical configurations, preserving positive gains and ties.}
\end{figure}

\subsection{Recovery Paths and Interference}

\paragraph{Available physical calibration.}
Figure~\ref{fig:eval-recovery-paths} reorganizes the existing path timings
used in the earlier reconstruction comparison. The configuration is Qwen2.5-7B-Instruct-1M,
two A100 PCIe 40\,GB GPUs with TP2, BF16, and batch size one. One fixed
synthetic token stream is evaluated at 8K, 64K, 256K, and approximately
1.01M tokens. HBM/DRAM reuse use three timing repetitions per input;
full Prefill and DPC shared-filesystem reuse use one. The endpoint is
the same input KV ready for continuation, including retrieval and
required computation. These are path measurements, not independent
workflow campaigns.

At the longest context, resident HBM, host DRAM, shared filesystem,
and full reconstruction take 1.3050, 2.6632, 10.1561, and 571.2127\,s.
The large reconstruction gap motivates retained state; the smaller
HBM--DRAM gap is the quantity a Peer tier could potentially reduce.
This measurement does not include a Peer transfer, background serving,
concurrent recovery, or fault detection. The filesystem server cache is
uncontrolled, so this is neither cold-SSD latency nor a bandwidth claim.

\begin{figure}[tb]
 \centering
 \includegraphics[width=\columnwidth]{figures/evaluation/eval-recovery-paths.pdf}
 \caption{Measured KV-path latency on one fixed input stream. HBM and
 DRAM show medians and observed ranges of three timing repetitions;
 reconstruction and shared-filesystem points each have one observation.
 HBM is resident reuse, and SharedFS is DPC shared storage. These are
 calibration measurements, not fault-injected system results.}
 \label{fig:eval-recovery-paths}
 \Description{Four measured recovery-related path latencies across four
 context lengths, with observed timing ranges only where repetitions exist.}
\end{figure}

\paragraph{Required causal decomposition.}
Log detection, fencing/rerouting, admission, directory lookup, storage
read, network transfer, H2D, required replay, request resume, and DAG
completion. Trace dependency edges and identify the realized critical
path; adding overlapping component durations would overcount recovery.
Hold demand and coverage fixed to compare Peer off/on with overlap
off/on. Separately measure total bytes, unused prefetch, GPU work,
temporary copies, and recovery-pool busy time. A faster request does not
necessarily reduce join delay or total recovery work.

Inject bursts affecting 1/2/4/8 workflows with 1/2/4 surviving targets
under the same foreground arrivals. Plot queue occupancy and
interrupted/uninterrupted WCT, TTFT, and TPOT over time, with an isolated
resource control. Shared recovery queues must be identified before
claiming collateral delay. This follows the cohort separation in
LUMEN~\cite{cao2026lumen} while retaining WCT as AegisKV's primary outcome.

\subsection{Mechanism Contributions and Credit Cost}
\label{sec:eval-rental}

Use matched interventions rather than removing several mechanisms at
once. The main ablations are: zero-slack fallback versus predicted slack;
reuse-only versus recovery-only versus combined priority; per-Agent
shares versus creator-only shares with physical deduplication unchanged;
causal-visibility normalization versus
raw selection counts; fixed-slot Slack versus FCFS, recency, and
recovery-cost admission; and fixed Peer capacity versus guarded AIMD.
For lower tiers, compare the current lazy-credit policy with same-value
direct revaluation, ordinary cost/size GDS or Landlord, no actual-use
renewal, and no pressure aging. Preserve the same ownership, leases,
arrival/fault opportunities, value estimates, and physical budgets.

Lazy versus eager refresh needs both an equal-CPU-budget comparison and
an unbounded eager reference. Report decision time, candidates examined,
scan age, stale rankings, differing victims, movement bytes, and paired
WCT. Fewer scalar checks alone do not establish a faster controller.
Five factorial experiments expose interactions: Slack $\times$ shared-Agent
shares, normalized demand $\times$ restoration pipelining, renewable credit
$\times$ lazy refresh, load-aware routing $\times$ shared-link contention,
and Peer residency $\times$ pipelining. Each uses all four cells with
the other features fixed. Campaign-level differences of differences
distinguish a combined effect from two separately observed improvements.

\paragraph{Available adverse result.}
The historical additive rental diagnostic has 512 decisions and the
same 396 value checks for direct revaluation and the original trial-rent
variant. Their miss-plus-movement costs are 462.96 and 506.39 model units:
trial rent is 9.38\% worse. Figure~\ref{fig:eval-rental-diagnostic} also
retains the later repaired variant. A diagnostic repair does not validate
the current per-Agent lazy-credit controller. This counterexample makes
the same-value direct baseline necessary; it is not evidence that adding
rent improves WCT or reduces CPU cost.

\begin{figure}[tb]
 \centering
 \includegraphics[width=\columnwidth]{figures/evaluation/eval-rental-diagnostic.pdf}
 \caption{Historical rental diagnostic. Miss and movement costs are
 retained for direct revaluation, the original trial-rent variant, and
 its repaired variant. Lower is better. The original negative result is
 preserved; model cost and value-check counts are not physical WCT or
 controller CPU measurements.}
 \label{fig:eval-rental-diagnostic}
 \Description{Historical miss-plus-movement cost comparison, preserving
 the rental regression and the remaining gap after a diagnostic repair.}
\end{figure}

\subsection{Prediction Quality and Sparse Demand}
\label{sec:eval-prediction}

\paragraph{Remaining time.}
Figure~\ref{fig:eval-tracker-errors} uses 1,443 eligible second-request
agent snapshots from 170 workflows, with one snapshot per Agent and
five workflow-disjoint folds. The target is whole-Agent remaining wall
time, not the residual segment duration needed by the current Tracker.
Mean absolute errors are 438.6\,s for a training-fold median, 406.7\,s
for basic GBDT, and 407.1\,s for GBDT with history. Basic GBDT's p95
absolute error is nevertheless worse than the median baseline
(1113.8 versus 1036.1\,s). The plot preserves
the complete error tail, including errors above 85,000\,s. A modest
average reduction does not establish trustworthy slack, and these
labels cannot validate segment-level predictions.

The current Tracker requires new causal segment labels, workflow-disjoint
training/calibration/test sets, and later-workload validation. Measure
remaining-time error, false positive slack on truly critical branches,
critical-path ranking, fallback frequency, and predictor/update cost.
Evaluate zero-slack and oracle-duration references. Bootstrap whole
workflows for prediction diagnostics; do not infer an error CDF from
aggregate WAPE or feature correlations. Study stale profiles, model
shifts, missing DAG edges, and delayed tools by recomputing decisions
against actual outcomes, rather than adding decorative error bars.

\begin{figure}[tb]
 \centering
 \includegraphics[width=\columnwidth]{figures/evaluation/eval-tracker-errors.pdf}
 \caption{Held-out absolute error of whole-Agent remaining-time
 predictions. The ECDF contains all 1,443 Agent snapshots from 170
 workflows; folds separate workflows. It is descriptive and has no
 invented confidence band. Whole-Agent accuracy does not validate
 residual-segment slack.}
 \label{fig:eval-tracker-errors}
 \Description{Complete held-out absolute-error ECDFs for median, basic
 GBDT, and history-augmented GBDT, including the long error tail.}
\end{figure}

\paragraph{Native access and horizon sensitivity.}
Figure~\ref{fig:eval-sparse-demand} compares historical frequency,
recency, and an offline demand oracle at the same within-model byte
budget. The captures use 32,768 historical tokens, a preceding
512-Query ranking window, and the following 64 Queries for evaluation.
Each object is a layer's 64-token block including all KV heads. The
budget is 205 of 512 historical blocks per layer, approximately 40\%.
Coverage is the fraction of distinct future-needed historical bytes
already selected, with equal document weighting and three documents
per model. It is neither attention mass nor a GPU hit-rate measurement.

Historical frequency covers 77.82\%, 53.57\%, and 43.02\% for
NOSA-3B, MiniCPM4-8B, and the experimental SeerAttention/Qwen3 path;
same-budget oracle coverage is 93.75\%, 58.98\%, and 43.28\%.
The small Seer oracle gap limits the value of a better ranker at this
budget. Extending the future window to 512 Queries lowers frequency
coverage to 58.90\%, 44.08\%, and 40.47\%; the complete horizon sweep
is retained in the artifact. These historical counts include native
head/group selection events. They are a baseline for, not a measurement
of, the current deduplicated, visibility-normalized probability estimate.

The required current-policy experiment compares raw counts, normalized
counts, recency, and random selection using strictly preceding windows.
An equal-budget offline oracle may use actual future demand and is
reported only as a clairvoyant reference. Sweep history length, future horizon,
budget, and unseen documents/workflows; account for selector state and
all physical objects. Freeze $L/K/Z$ on calibration workloads and test
them on held-out workloads. Finally measure exposed wait and WCT under
misses and foreground contention: high coverage alone does not prove
net Peer benefit.

\begin{figure}[tb]
 \centering
 \includegraphics[width=\columnwidth]{figures/evaluation/eval-sparse-demand.pdf}
 \caption{Native sparse-access diagnostic at an equal 205-of-512-block
 budget per layer. Historical frequency and recency select for the
 following 64 Queries; the offline oracle sees their realized demand.
 Points and summaries retain the
 three-document variation; this is observed demand coverage, not a
 confidence interval or failure-recovery speedup.}
 \label{fig:eval-sparse-demand}
 \Description{Same-budget native demand coverage for three models and
 three placement references, with the individual document variation retained.}
\end{figure}

\subsection{Healthy Service and Sustained Failures}

\paragraph{Normal-service cost.}
Run the complete original workloads under $F0$, including first calls,
tool waits, and follow-up requests. Measure WCT, TTFT, TPOT, SLO
goodput, GPU-seconds, controller CPU/memory, checkpoint bytes, and
per-tier request and token hits. Hold absolute arrival traces fixed at
30/60/85\% of the reference capacity; do not normalize each policy to
its own throughput. Record execution-HBM loss and migration contention
when Peer grows. Prefix overlap in a trace is not a runtime cache hit.
Report conditional recovery gains separately from any amortized benefit
at a declared failure frequency. A model queueing delay is not measured
TTFT, and a failure-heavy stress test does not establish production net
benefit.

\paragraph{Why averages are insufficient.}
The historical LANL-timing model~\cite{lanlFailureData,schroeder2006failures}
retains 21,738 cleaned node-outage
intervals from 22 systems. Each start triggers an assumed four-branch
recovery in a per-system FIFO spare pool. Figure~\ref{fig:eval-trace-tails}
shows a lower event-weighted mean for workflow-value allocation
(9.4673 versus 10.4533\,s) but worse p95 (40.0292 versus 37.0375\,s)
and p99 (232.1125 versus 204.1208\,s). Its longer pool occupancy can
hurt subsequent recoveries even when the current join finishes sooner.
These are modeled outcomes driven by recorded timing, not observed
LANL agent workflows or measured AegisKV failures.

\begin{figure}[tb]
 \centering
 \includegraphics[width=\columnwidth]{figures/evaluation/eval-trace-tails.pdf}
 \caption{Historical trace-driven recovery model: mean and tail failure
 penalties under the same recorded event timing. Workflow-value
 placement reduces the mean while worsening p95 and p99. The outcome
 combines an assumed workload and spare-pool model with recorded outage
 starts; it is not measured fleet availability.}
 \label{fig:eval-trace-tails}
 \Description{Mean and tail recovery penalties in a historical FIFO
 model, retaining the adverse p95 and p99 comparison.}
\end{figure}

For continuous replay, preserve failure clusters, downtime, surviving
capacity, and non-overlapping calendar blocks; report event-weighted
and system/window-weighted summaries. The available LANL redistribution
has 23,739 rows versus 23,741 documented by its original source, an
unresolved provenance difference. Reflex/ByteDance
alerts~\cite{ding2026reflex,reflexDataset} are shifted
within each job and lack reliable down/up pairs, so they cannot define
a global failure calendar. Modern node-fault traces such as
InfiniteHBD~\cite{shou2025infinitehbd} also require their original
node-level semantics and training-cluster scope. Trace provenance and
explicit resource mapping precede any claim about serving availability.

\subsection{Operating Range and Controller Scaling}

Capacity experiments ask when protection loses its advantage. Independent
Peer, DRAM, and storage sweeps include Peer-off, full destinations, and
last-copy pressure under equal physical budgets. Achieved coverage and
commit lag distinguish insufficient capacity from an overloaded copying
path: the uncommitted tail $N$ may exceed the checkpoint interval $Z$.
WCT is therefore interpreted alongside write backlog, migration traffic,
unused prefetch, and fallback frequency. Separate storage, NIC, and H2D
bottlenecks reveal whether a gain depends on one favorable transport ratio.

Scaling experiments separate replica count, TP degree, and failure domains
from controller state size. Increasing DAG size, shared-owner fan-out,
live objects, and candidate budget exposes the tradeoff between decision
quality and update cost. Actual decision-time tails, CPU use, metadata,
and foreground stalls compare bounded lazy refresh with equal-budget
and unbounded eager controls. Saturation and unsupported cells remain
visible; multiplying operation counts by an assumed node cost cannot
substitute for controller timing.

Robustness experiments test whether fallback and capacity control remain
useful when duration errors correlate across a DAG, profiles become stale,
or tools and arrivals change abruptly. Step, ramp, and burst loads expose
Adapter overshoot and settling time. The key capacity outcome is actually
freed bytes: lowering a quota target does not release memory until the
Peer-to-DRAM-to-storage cascade and reader drain finish.

\subsection{Recovery Correctness and Reproducibility}
\label{sec:eval-plan}

Correctness is a prerequisite for interpreting any latency result. The
state-transition suite checks exclusive coverage, commit-before-redirect,
and reader-drain-before-free under stale epochs, duplicate events, and
incomplete restore state. Sole-owner loss must lead to a declared
older-checkpoint/replay fallback. Pressure and tier-round-trip cases
check physical occupancy, retained paid rent, and cancellation; Peer
cases check the specified victim order, tie rules, and physical leases.
The artifact enumerates the injection points and seeded interleavings.

Faults bracketing copy, commit, ownership switch, and acknowledgement
test accepted-token durability and complete compatible restore state.
Next-output equivalence uses a declared numerical tolerance, while
acknowledged output must never be duplicated. Application tool
idempotency is a separate requirement. Transition coverage and failed
randomized schedules accompany the results; earlier replicated-Peer
examples cannot establish these invariants for exclusive ownership.

Freeze source/model/backend versions, workload and fault IDs, physical
budgets, and calibration choices before collecting primary results.
Randomize policy order and start with ten paired campaigns; increase
repetition based on pilot variance and precision rather than the desired
sign of an effect. Aggregate by campaign and publish bootstrap method,
seed, and intervals. Deterministic model scenarios have no experimental
confidence intervals; the microbenchmark ranges describe repeated
timings of one input, not independent workflow uncertainty. Preserve
timeouts in the completion denominator and do not invent finish times.

The artifact contains the complete experiment matrix, per-figure source
hashes, long-form plotting tables, native-size vector exports, and
generation commands. All unexecuted system cells remain explicitly
planned. The execution order is correctness and backend calibration,
paired single-failure outcomes, healthy overhead and decisive ablations,
then continuous trace replay and scale/pressure tests. Attractive plots
and a complete protocol do not close the current gap in end-to-end
implementation evidence.


\section{Related Work}
\label{sec:related}

\paragraph{Fault-tolerant LLM serving.}
D\'ej\`aVu streams and replicates KV~\cite{strati2024dejavu}; GhostServe
protects it with host-memory erasure-coded parity~\cite{jayakody2026ghostserve}.
RaidServe addresses TP resilience and imbalance through cyclic placement,
hybrid attention and recovery routing~\cite{xu2026raidserve}. LUMEN
coordinates checkpoints, request redistribution and capacity restoration
under reduced load~\cite{cao2026lumen}. Placement comparisons must use the same strong recovery scheduler.

\paragraph{Hierarchical context caching.}
Mooncake pools KV and schedules around locality~\cite{qin2025mooncake};
its SSD offload spills state under memory pressure~\cite{mooncake2026ssdoffload}.
CachedAttention uses hierarchical caching, layer-wise preload and
asynchronous saving~\cite{gao2024cachedattention}. Strata addresses
fragmented I/O, delay hits and cache-aware scheduling~\cite{xie2026strata}.
HCache restores intermediate activations while balancing I/O and
computation~\cite{gao2025hcache}. Our proposed distinction is common-failure workflow value
with safe demotion and explicit final-copy fallback.

\paragraph{Workflow-aware serving and prediction.}
Parrot optimizes application-level execution~\cite{lin2024parrot}; KVFlow
uses future invocation distance for prefix eviction~\cite{pan2025kvflow}.
Pythia and Cortex exploit workflow semantics~\cite{yu2026pythia,pagonas2025cortex}.
TOPAS schedules retention and execution using remaining paths, reuse and
movement cost~\cite{ni2026topas}; adding a critical-path score to caching
is thus insufficient novelty. The intended distinction is recovery after
a common worker failure with stale checkpoints and finite-tier transitions.
Heddle predicts trajectories~\cite{zhang2026heddle}, which does not establish
AegisKV's remaining-wall-time accuracy.

\paragraph{Placement mechanisms.}
Critical-path analysis~\cite{kelley1959cpm} and cost/size-based caching~\cite{young2002landlord}
provide building blocks. AegisKV's values depend on the DAG, shared queues,
and object locations, so Landlord's competitive bound does not transfer. The
planned small-instance feasible oracle will quantify approximation loss. AegisKV
uses lower-tier renewable credit with lazy refresh and factorized request/block
priority, alongside fixed-quota Peer admission. Independent ablations against
direct revaluation must establish the benefit of rent, actual-use renewal
and value refresh, alongside sparse ranking and overlap, under matched
resource and control-work budgets.

\section{Conclusion}
\label{sec:conclusion}

AegisKV combines normal reuse, workflow urgency and block demand to manage
finite Peer HBM, DRAM and shared storage. Exclusive ownership distributes
complete checkpoints across these tiers; fixed-quota hot sets accelerate
recovery. The online controller combines lazy lower-tier credit with greedy
Peer admission, distinguishing migration from final-copy deletion and
preserving committed coverage through each transition. Analytical
scenarios identify potential benefits and limits. A historical trace-driven
recovery model shows that average improvements can coexist with tail
regressions; isolated rental diagnostics reveal costs that require fresh
validation. These
prior model results do not validate the exclusive three-tier controller. The detailed
protocol and single-column figures support the physical correctness and performance
experiments still needed to establish system-level gains.

\paragraph{AI assistance.}
Generative AI assisted drafting, experiment planning, and artifact code;
authors retain responsibility for verification.

\bibliographystyle{ACM-Reference-Format}
\bibliography{references}

% BEGIN MOTIVATION MEASUREMENT APPENDIX
\appendix
\section{Motivation Measurement Protocol}
\label{sec:motivation-protocol}

\paragraph{Scope and configuration.}
Figures~\ref{fig:motivation-recompute} and
\ref{fig:motivation}(\subref{fig:motivation-recovery-exposure}) use the same pinned
Qwen2.5-7B-Instruct-1M weights, SGLang 0.5.10rc0, PyTorch 2.9.1/CUDA 12.8,
TP2, BF16 weights/KV, and batch one on two A100 PCIe 40\,GB GPUs.
They use separate allocations of the same GPU type, with one fixed
repeated synthetic text tokenized by the model tokenizer. The traces
supply workflow positions and length proxies, not benchmark prompts.
Our adapter uses the pinned MInference vertical/slash configuration and
DCA position handling, with sparse attention above 32,768 tokens,
including the final one-token EXTEND. The SM80 sparse-kernel build and
model/attention sources are archived with hashes. GQA K/V are shared;
prefix execution groups at most two query heads, and a single query
batches all heads. These are instrumented worker measurements, rather
than throughput measurements of an unmodified serving frontend.

\paragraph{Common endpoint.}
For $C=M+1$, full prefill computes $M$ tokens from empty KV in chunks of
at most 8,192, then executes one token separately. Reuse consumes exactly
the same sparse-generated $M$-token state, including partial pages, and
executes that final token. Synchronized wall time is reduced by maximum
across the TP ranks and ends at model-body completion with all $C$ tokens'
KV resident, before the output vocabulary projection. Model loading,
backup creation, and correctness checks are outside the timer. Detection,
rerouting, serving-queue delay and an injected failure are absent.

\paragraph{Restore paths.}
HBM begins with resident KV. DRAM copies all historical layers from
pinned host buffers before computation. Shared storage opens, reads and
closes per-rank KV files on the compute node's DPC shared filesystem,
then performs the same H2D copy and computation. These stages do not
overlap; there is no compression or selective loading. The 1,010,000-token
case restores 57,917,382,656 logical history bytes across both ranks;
files additionally contain recorded alignment padding. Direct I/O and
page-residency checks exclude Linux file-page hits, but DPC server and
internal caches are uncontrolled. The curve is a shared-storage path
measurement, not cold-device latency or an identified SSD measurement.
KV files are removed after measurement.

\paragraph{Timing samples.}
Figure~\ref{fig:motivation-recompute} uses one
allocation for 8,192, 65,536, 262,144 and 1,010,000 tokens. HBM/DRAM have
one warmup and three timings per length; full prefill and shared storage
have one timing each. Short forwards and operator checks precede timing;
64K also runs a complete original all-head sparse reference. There is
no extra complete-prefill warmup at 256K/1.01M or storage-read warmup.
Consequently, these full-prefill/storage points describe individual
runs, not stable latency distributions.
Figure~\ref{fig:motivation}(\subref{fig:motivation-recovery-exposure}) uses two allocations at
$C=100{,}865$, $105{,}921$, $108{,}865$ and $113{,}665$, with one complete
warmup and three full-prefill timings at each length. The initial allocation completed three lengths; the last length was
resumed after an early termination signal. Its incomplete timings are
retained in the artifact but excluded as a whole. The figure reports
medians and min--max ranges. Repetitions are timing samples of the same input,
not independent workflows or confidence intervals. Original timings and
all samples are retained in the artifact.

\paragraph{Validation and workflow mapping.}
Complete-history hashes, final hidden states, all-layer new KV and greedy
outputs are checked across reconstruction/reuse paths. Repeated full
prefills additionally check output agreement, and CUDA-event timelines
check serialized copies. These checks validate matched sparse-state
execution, not equivalence to dense attention or long-context task quality.
For Figure~\ref{fig:motivation}(\subref{fig:motivation-recovery-exposure}), the four history lengths
are the active requests' trace hash-block counts multiplied by 64.
They are size proxies, not exact original-model token counts or observed
KV frontiers at failure. The example was selected for four parallel
branches and a gap near 100\,s; it is not its workflow's maximum-gap group.
Each counterfactual separately inserts a measured full-prefill duration
as an additional branch pause at 29.5\,s after the fork. The all-branch
join, other execution and queues stay fixed. Thus only reconstruction
cost is measured: the plotted WCT increment is a conditional calculation,
not observed fault recovery in the Claude Code workflow.
% END MOTIVATION MEASUREMENT APPENDIX
\end{document}
